\documentclass[10pt,a4paper]{amsart}
\usepackage{amsmath,amssymb,amsthm,graphicx}
% For commutative diagrams we use Xy-pic
\usepackage[all]{xy}

% We use 2cell for 2-commutative diagrams.
\xyoption{2cell}
% Theorem environments.
%
\theoremstyle{plain}
\newtheorem{theorem}[subsection]{Theorem}
\newtheorem{proposition}[subsection]{Proposition}
\newtheorem{lemma}[subsection]{Lemma}

\theoremstyle{definition}
\newtheorem{definition}[subsection]{Definition}
\newtheorem{example}[subsection]{Example}
\newtheorem{exercise}[subsection]{Exercise}
\newtheorem{situation}[subsection]{Situation}

\theoremstyle{remark}
\newtheorem{remark}[subsection]{Remark}
\newtheorem{remarks}[subsection]{Remarks}

\numberwithin{equation}{subsection}

% Macros
%
\def\lim{\mathop{\mathrm{lim}}\nolimits}
\def\colim{\mathop{\mathrm{colim}}\nolimits}
\def\Spec{\mathop{\mathrm{Spec}}}
\def\Hom{\mathop{\mathrm{Hom}}\nolimits}
\def\Ext{\mathop{\mathrm{Ext}}\nolimits}
\def\SheafHom{\mathop{\mathcal{H}\!\mathit{om}}\nolimits}
\def\SheafExt{\mathop{\mathcal{E}\!\mathit{xt}}\nolimits}
\def\Sch{\mathit{Sch}}
\def\Mor{\mathop{\mathrm{Mor}}\nolimits}
\def\Ob{\mathop{\mathrm{Ob}}\nolimits}
\def\Sh{\mathop{\mathit{Sh}}\nolimits}
\def\NL{\mathop{N\!L}\nolimits}
\def\CH{\mathop{\mathrm{CH}}\nolimits}
\def\proetale{{pro\text{-}\acute{e}tale}}
\def\etale{{\acute{e}tale}}
\def\QCoh{\mathit{QCoh}}
\def\Ker{\mathop{\mathrm{Ker}}}
\def\Im{\mathop{\mathrm{Im}}}
\def\Coker{\mathop{\mathrm{Coker}}}
\def\Coim{\mathop{\mathrm{Coim}}}

% Boxtimes
%
\DeclareMathSymbol{\boxtimes}{\mathbin}{AMSa}{"02}

%
% Macros for moduli stacks/spaces
%
\def\QCohstack{\mathcal{QC}\!\mathit{oh}}
\def\Cohstack{\mathcal{C}\!\mathit{oh}}
\def\Spacesstack{\mathcal{S}\!\mathit{paces}}
\def\Quotfunctor{\mathrm{Quot}}
\def\Hilbfunctor{\mathrm{Hilb}}
\def\Curvesstack{\mathcal{C}\!\mathit{urves}}
\def\Polarizedstack{\mathcal{P}\!\mathit{olarized}}
\def\Complexesstack{\mathcal{C}\!\mathit{omplexes}}
% \Pic is the operator that assigns to X its picard group, usage \Pic(X)
% \Picardstack_{X/B} denotes the Picard stack of X over B
% \Picardfunctor_{X/B} denotes the Picard functor of X over B
\def\Pic{\mathop{\mathrm{Pic}}\nolimits}
\def\Picardstack{\mathcal{P}\!\mathit{ic}}
\def\Picardfunctor{\mathrm{Pic}}
\def\Deformationcategory{\mathcal{D}\!\mathit{ef}}
\usepackage{cleveref}
\crefname{originaldefinition}{original definition}{original definitions}
\Crefname{originaldefinition}{Original definition}{Original definitions}
\makeatletter
\newlabel{situation-setup}{{02QL}{source}{}{sourcepointer-1}{}}
\newlabel{situation-setup@cref}{{[originaldefinition][0][]02QL}{[1][0][]source}{}{}{}}
\makeatother
\crefname{originaldefinition}{original definition}{original definitions}
\Crefname{originaldefinition}{Original definition}{Original definitions}
\makeatletter
\newlabel{divisors-definition-regular-meromorphic-ideal-denominators}{{02P1}{source}{}{sourcepointer-2}{}}
\newlabel{divisors-definition-regular-meromorphic-ideal-denominators@cref}{{[originaldefinition][0][]02P1}{[1][0][]source}{}{}{}}
\makeatother
\crefname{originallemma}{original lemma}{original lemmas}
\Crefname{originallemma}{Original lemma}{Original lemmas}
\makeatletter
\newlabel{divisors-lemma-make-maps-regular-section}{{02P2}{source}{}{sourcepointer-3}{}}
\newlabel{divisors-lemma-make-maps-regular-section@cref}{{[originallemma][0][]02P2}{[1][0][]source}{}{}{}}
\makeatother
\crefname{originallemma}{original lemma}{original lemmas}
\Crefname{originallemma}{Original lemma}{Original lemmas}
\makeatletter
\newlabel{divisors-lemma-pullback-meromorphic-sections-defined}{{02OU}{source}{}{sourcepointer-4}{}}
\newlabel{divisors-lemma-pullback-meromorphic-sections-defined@cref}{{[originallemma][0][]02OU}{[1][0][]source}{}{}{}}
\makeatother
\crefname{originallemma}{original lemma}{original lemmas}
\Crefname{originallemma}{Original lemma}{Original lemmas}
\makeatletter
\newlabel{divisors-lemma-meromorphic-sections-pullback}{{02OY}{source}{}{sourcepointer-5}{}}
\newlabel{divisors-lemma-meromorphic-sections-pullback@cref}{{[originallemma][0][]02OY}{[1][0][]source}{}{}{}}
\makeatother
\crefname{originallemma}{original lemma}{original lemmas}
\Crefname{originallemma}{Original lemma}{Original lemmas}
\makeatletter
\newlabel{lemma-no-embedded-points-modules}{{02QH}{source}{}{sourcepointer-6}{}}
\newlabel{lemma-no-embedded-points-modules@cref}{{[originallemma][0][]02QH}{[1][0][]source}{}{}{}}
\makeatother
\title{The ideal of denominators of a regular meromorphic section: two short exact sequences with support condition and a cycle-theoretic difference formula}
\author{The Stacks Project Authors}
\date{stacks-project:chow.tex:local-snapshot-20261012; GNU FDL 1.2 or later}
\begin{document}
\maketitle
\noindent\textit{Editorial note.} All delivered mathematical text is the author's own, copied from the pinned Stacks Project chapter file chow.tex; no statement or proof step was rewritten, reordered or removed.\par\medskip
\noindent\textit{Editorial note.} Preamble provenance: the pinned chapter file opens with an input of the shared preamble and carries no preamble text of its own, so the project's single Stacks-wide preamble staged at sources/stacks-chow/preamble.tex is excerpted verbatim as two spans, lines 14-18 (Xy-pic setup) and lines 154-217 (the native theorem declarations and mathematical macros). Nothing was added, removed or redefined; the wrapper class and the amsmath/amssymb/amsthm/graphicx base packages are supplied by the assembler.\par\medskip
\noindent\textit{Editorial note.} The counted statement and proof contain no cross reference, no citation and no included file, so the delivered body is self contained.\par\medskip
\noindent\textit{Editorial note.} Canonical location of the counted proposition: Stacks tag 02SW (sources/stacks-chow/pinned-original-source.tex line 18649).\par\medskip
\noindent\textit{Allowed standard prerequisites (editorial).}\begin{itemize}
\item The ideal of denominators of a regular meromorphic section and the associated injections. Divisors, Definition labelled divisors-definition-regular-meromorphic-ideal-denominators (Stacks tag 02P1) and Lemma labelled divisors-lemma-make-maps-regular-section (Stacks tag 02P2), invoked at statement line 18667 and proof lines 18715-18719. Supply the ideal sheaf I of denominators and the injective maps 1 : I F -> F and s : I F -> F tensor L whose cokernels are supported on a nowhere dense closed subset; the identification of the quotient sheaves and the cycle computation are inside the counted span.
\item Pullback of meromorphic sections along the components and regularity of the restrictions. Divisors, Lemmas labelled divisors-lemma-pullback-meromorphic-sections-defined (Stacks tag 02OU) and divisors-lemma-meromorphic-sections-pullback (Stacks tag 02OY), invoked at proof lines 18722-18724. Supply that the restrictions s\_i to the components X\_i of delta-dimension k+1 are defined and regular meromorphic for L|\_\{X\_i\}.
\item The module-level comparison of the two lengths with the divisor sum. Same chapter, Lemma labelled lemma-no-embedded-points-modules (source line 18594), invoked at proof line 18752. Supplies the numerical identity that makes the coefficient of [Z] on the two sides agree; the reduction to it, including the rewriting s = (a/b) s\_xi and the identification of the coefficient formulas, is carried out in the counted span.
\item The base situation (S, delta) for operational Chow groups and cycle classes. Same chapter, Situation labelled situation-setup (source line 1372), invoked in the statement. Fixes the delta-dimension conventions, the cycle group Z\_k(X) and the pullback and cap operations used in the statement.
\end{itemize}\medskip
\ifdefined\hypertarget\hypertarget{sourcepointer-1}{}\fi
\noindent\textit{Source reference.} Original definition 02QL: sources/stacks-chow/pinned-original-source.tex, source lines 1372--1372. Same-chapter pointer into the author environment situation at source line 1372; invoked at proof line None of the counted proof. It is retained only as a navigation pointer; every load-bearing step is inside the delivered span.\par\medskip
\ifdefined\hypertarget\hypertarget{sourcepointer-2}{}\fi
\noindent\textit{Source reference.} Original definition 02P1: sources/stacks-chow/tags.txt, source lines 3536--3536. Cross-chapter standard prerequisite located in the project tag list staged in the local snapshot; invoked at proof line None of the counted proof. It is retained only as a navigation pointer; every load-bearing step is inside the delivered span.\par\medskip
\ifdefined\hypertarget\hypertarget{sourcepointer-3}{}\fi
\noindent\textit{Source reference.} Original lemma 02P2: sources/stacks-chow/tags.txt, source lines 3537--3537. Cross-chapter standard prerequisite located in the project tag list staged in the local snapshot; invoked at proof line 18715 of the counted proof. It is retained only as a navigation pointer; every load-bearing step is inside the delivered span.\par\medskip
\ifdefined\hypertarget\hypertarget{sourcepointer-4}{}\fi
\noindent\textit{Source reference.} Original lemma 02OU: sources/stacks-chow/tags.txt, source lines 3529--3529. Cross-chapter standard prerequisite located in the project tag list staged in the local snapshot; invoked at proof line 18722 of the counted proof. It is retained only as a navigation pointer; every load-bearing step is inside the delivered span.\par\medskip
\ifdefined\hypertarget\hypertarget{sourcepointer-5}{}\fi
\noindent\textit{Source reference.} Original lemma 02OY: sources/stacks-chow/tags.txt, source lines 3533--3533. Cross-chapter standard prerequisite located in the project tag list staged in the local snapshot; invoked at proof line 18723 of the counted proof. It is retained only as a navigation pointer; every load-bearing step is inside the delivered span.\par\medskip
\ifdefined\hypertarget\hypertarget{sourcepointer-6}{}\fi
\noindent\textit{Source reference.} Original lemma 02QH: sources/stacks-chow/pinned-original-source.tex, source lines 18594--18594. Same-chapter pointer into the author environment lemma at source line 18594; invoked at proof line 18752 of the counted proof. It is retained only as a navigation pointer; every load-bearing step is inside the delivered span.\par\medskip
\begin{lemma}
\label{lemma-no-embedded-points}
Let $(S, \delta)$ be as in Situation \ref{situation-setup}.
Let $X$ be locally of finite type over $S$.
Let $\mathcal{L}$ be an invertible $\mathcal{O}_X$-module.
Let $\mathcal{F}$ be a coherent $\mathcal{O}_X$-module.
Let $s \in \Gamma(X, \mathcal{K}_X(\mathcal{L}))$ be a
meromorphic section of $\mathcal{L}$.
Assume
\begin{enumerate}
\item $\dim_\delta(X) \leq k + 1$,
\item $X$ has no embedded points,
\item $\mathcal{F}$ has no embedded associated points,
\item the support of $\mathcal{F}$ is $X$, and
\item the section $s$ is regular meromorphic.
\end{enumerate}
In this situation let $\mathcal{I} \subset \mathcal{O}_X$
be the ideal of denominators of $s$, see
Divisors,
Definition \ref{divisors-definition-regular-meromorphic-ideal-denominators}.
Then we have the following:
\begin{enumerate}
\item there are short exact sequences
$$
\begin{matrix}
0 &
\to &
\mathcal{I}\mathcal{F} &
\xrightarrow{1} &
\mathcal{F} &
\to &
\mathcal{Q}_1 &
\to &
0 \\
0 &
\to &
\mathcal{I}\mathcal{F} &
\xrightarrow{s} &
\mathcal{F} \otimes_{\mathcal{O}_X} \mathcal{L} &
\to &
\mathcal{Q}_2 &
\to &
0
\end{matrix}
$$
\item the coherent sheaves $\mathcal{Q}_1$, $\mathcal{Q}_2$
are supported in $\delta$-dimension $\leq k$,
\item the section $s$ restricts to a regular meromorphic
section $s_i$ on every irreducible component $X_i$ of
$X$ of $\delta$-dimension $k + 1$, and
\item writing $[\mathcal{F}]_{k + 1} = \sum m_i[X_i]$ we have
$$
[\mathcal{Q}_2]_k - [\mathcal{Q}_1]_k
=
\sum m_i(X_i \to X)_*\text{div}_{\mathcal{L}|_{X_i}}(s_i)
$$
in $Z_k(X)$, in particular
$$
[\mathcal{Q}_2]_k - [\mathcal{Q}_1]_k
=
c_1(\mathcal{L}) \cap [\mathcal{F}]_{k + 1}
$$
in $\CH_k(X)$.
\end{enumerate}
\end{lemma}

\begin{proof}
Recall from Divisors, Lemma \ref{divisors-lemma-make-maps-regular-section}
the existence of injective maps
$1 : \mathcal{I}\mathcal{F} \to \mathcal{F}$ and
$s : \mathcal{I}\mathcal{F} \to \mathcal{F} \otimes_{\mathcal{O}_X}\mathcal{L}$
whose cokernels are supported on a closed nowhere dense subsets $T$.
Denote $\mathcal{Q}_i$ there cokernels as in the lemma.
We conclude that $\dim_\delta(\text{Supp}(\mathcal{Q}_i)) \leq k$.
By Divisors, Lemmas \ref{divisors-lemma-pullback-meromorphic-sections-defined}
and \ref{divisors-lemma-meromorphic-sections-pullback} the pullbacks $s_i$
are defined and are regular meromorphic sections for $\mathcal{L}|_{X_i}$.
The equality of cycles in (4) implies the equality of cycle classes
in (4). Hence the only remaining thing to show is that
$$
[\mathcal{Q}_2]_k - [\mathcal{Q}_1]_k
=
\sum m_i(X_i \to X)_*\text{div}_{\mathcal{L}|_{X_i}}(s_i)
$$
holds in $Z_k(X)$. To see this, let $Z \subset X$ be an integral closed
subscheme of $\delta$-dimension $k$. Let $\xi \in Z$ be the generic point.
Let $A = \mathcal{O}_{X, \xi}$ and $M = \mathcal{F}_\xi$.
Moreover, choose a generator $s_\xi \in \mathcal{L}_\xi$.
Then we can write $s = (a/b) s_\xi$ where $a, b \in A$ are
nonzerodivisors. In this case
$I = \mathcal{I}_\xi = \{x \in A \mid x(a/b) \in A\}$.
In this case the coefficient of $[Z]$ in the left hand side is
$$
\text{length}_A(M/(a/b)IM) - \text{length}_A(M/IM)
$$
and the coefficient of $[Z]$ in the right hand side
is
$$
\sum
\text{length}_{A_{\mathfrak q_i}}(M_{\mathfrak q_i})
\text{ord}_{A/\mathfrak q_i}(a/b)
$$
where $\mathfrak q_1, \ldots, \mathfrak q_t$ are the minimal
primes of the $1$-dimensional local ring $A$. Hence the result
follows from Lemma \ref{lemma-no-embedded-points-modules}.
\end{proof}

\end{document}
